% Options for packages loaded elsewhere
\PassOptionsToPackage{unicode}{hyperref}
\PassOptionsToPackage{hyphens}{url}
\documentclass[
]{article}
\usepackage{xcolor}
\usepackage{amsmath,amssymb}
\setcounter{secnumdepth}{-\maxdimen} % remove section numbering
\usepackage{iftex}
\ifPDFTeX
  \usepackage[T1]{fontenc}
  \usepackage[utf8]{inputenc}
  \usepackage{textcomp} % provide euro and other symbols
\else % if luatex or xetex
  \usepackage{unicode-math} % this also loads fontspec
  \defaultfontfeatures{Scale=MatchLowercase}
  \defaultfontfeatures[\rmfamily]{Ligatures=TeX,Scale=1}
\fi
\usepackage{lmodern}
\ifPDFTeX\else
  % xetex/luatex font selection
\fi
% Use upquote if available, for straight quotes in verbatim environments
\IfFileExists{upquote.sty}{\usepackage{upquote}}{}
\IfFileExists{microtype.sty}{% use microtype if available
  \usepackage[]{microtype}
  \UseMicrotypeSet[protrusion]{basicmath} % disable protrusion for tt fonts
}{}
\makeatletter
\@ifundefined{KOMAClassName}{% if non-KOMA class
  \IfFileExists{parskip.sty}{%
    \usepackage{parskip}
  }{% else
    \setlength{\parindent}{0pt}
    \setlength{\parskip}{6pt plus 2pt minus 1pt}}
}{% if KOMA class
  \KOMAoptions{parskip=half}}
\makeatother
\usepackage{longtable,booktabs,array}
\newcounter{none} % for unnumbered tables
\usepackage{calc} % for calculating minipage widths
% Correct order of tables after \paragraph or \subparagraph
\usepackage{etoolbox}
\makeatletter
\patchcmd\longtable{\par}{\if@noskipsec\mbox{}\fi\par}{}{}
\makeatother
% Allow footnotes in longtable head/foot
\IfFileExists{footnotehyper.sty}{\usepackage{footnotehyper}}{\usepackage{footnote}}
\makesavenoteenv{longtable}
\setlength{\emergencystretch}{3em} % prevent overfull lines
\providecommand{\tightlist}{%
  \setlength{\itemsep}{0pt}\setlength{\parskip}{0pt}}
\usepackage{hyperxmp}
\usepackage{bookmark}
\IfFileExists{xurl.sty}{\usepackage{xurl}}{} % add URL line breaks if available
\urlstyle{same}
\hypersetup{
  pdftitle={Program--value separability: the structural precondition for compilation{,} caching{,} and dense journaling in a DSL runtime},
  pdfauthor={Alvaro Rivera},
  pdfkeywords={\xmpquote{program-value
separability}, \xmpquote{code-as-data}, \xmpquote{homoiconic
persistence}, \xmpquote{DSL runtime}, \xmpquote{dual
compilation}, \xmpquote{journal density}, \xmpquote{partial
evaluation}, \xmpquote{CQRS}, \xmpquote{actor model}, \xmpquote{event
sourcing}, \xmpquote{puppeteer framework}},
  hidelinks,
  pdfcreator={LaTeX via pandoc}}

\title{Program--value separability: the structural precondition for
compilation, caching, and dense journaling in a DSL runtime}
\author{Alvaro Rivera}
\date{2026-05-18}

\begin{document}
\maketitle
\begin{abstract}
This paper is a design theory contribution in the sense of Hevner,
March, Park, and Ram (2004): it introduces a structural construct ---
\emph{program-value separability} --- derives the runtime consequences
that follow from it, and presents one realization in production as
confirmation that the construct is realizable. Compilation of a
domain-specific language program to native code, caching of compiled
programs across invocations, and dense journaling of operations rather
than their effects are commonly treated as independent runtime
optimizations. This paper argues that they are not --- that all three
are downstream consequences of a single structural property of DSL
programs, which we name \emph{program--value separability}: the
syntactically decidable condition that a program declares its parameters
explicitly rather than embedding values in its body. We trace four
runtime faces (compilation, caching, dense journaling,
replication-bounded entropy) as projections of this single precondition,
characterize one concrete realization in a CQRS + Actor + Event Sourcing
runtime --- the same Puppeteer framework introduced in the prior paper
of this series --- and report empirical magnitudes from measurements
over two independent open-source DDD aggregates (dotnet/eShop's Order
and kgrzybek/modular-monolith's SubscriptionPayment): a compiled-
versus-interpreted speedup that scales with DSL-bound work (1.49× to
4.10×); a cold compile cost amortized within hundreds of invocations; a
journal storing 99.8\% of parametric workloads as compact action
references, 3.5× to 5.6× denser than the equivalent literal-script
storage. The principle generalizes beyond any one runtime: program-value
separability characterizes the structural condition any DSL runtime must
satisfy to admit compilation, caching, and dense persistence at all.
\end{abstract}

\section{Program--value separability}\label{programvalue-separability}

\subsection{TL;DR}\label{tldr}

\begin{quote}
Compilation, caching, and dense journaling are \textbf{not} features
layered atop interpretation. They \textbf{are} downstream consequences
of a single structural commitment: values live outside the script, not
embedded within it.

The transition often described as ``adding compilation'' inverts the
causal order. A DSL whose programs embed their values has no stable
identifier, no template to cache, no separable sub-program to compile
ahead of time --- interpretation is not a stylistic choice but a
structural necessity. Once values move outside the script, the program
becomes a function awaiting its arguments --- \emph{F(x₁, x₂, \ldots,
xₙ)}, as the runtime documents itself.

From that single act of separation, four runtime decisions cohere.
Scripts with externalized parameters are eligible for compilation to IL
via Expression trees, cache as reusable programs with action
identifiers, persist to the journal as compact action entries, and
replicate with entropy bounded by argument vectors. Scripts with
hardcoded values are not cached, run once, and persist as their literal
text. The four faces are not parallel design choices; they are
projections of a single precondition.

This precondition operates on the substrate characterized in
\href{01-anti-porosity.md}{the previous paper of this series} ---
Puppeteer's journal as a homoiconic representation, code and data
sharing one form. We adopt the colloquial label \emph{code-as-data} from
this point on; the formalism stands as established. The compactness of
these entries is not merely syntactic: each names an operation
substantially richer than its surface signature, with the asymmetry
scaling with the domain library's depth (§5.5).

\textbf{This is the principle of program--value separability: a DSL
program becomes identifiable, compilable, cacheable, and persistable
densely only when it is separable from its values. Externalized
parameters are the structural precondition under which compilation,
caching, and dense journaling become possible at all.}
\end{quote}

\begin{center}\rule{0.5\linewidth}{0.5pt}\end{center}

\subsection{Claims this paper makes}\label{claims-this-paper-makes}

\begin{enumerate}
\def\labelenumi{\arabic{enumi}.}
\item
  \textbf{Causal inversion: program--value separability is a necessary
  condition.} Compilation, caching, and dense journaling in the runtime
  described here are not optimizations layered atop interpretation. They
  are downstream consequences of \emph{program--value separability} ---
  the structural property that values live outside the script body, not
  embedded within it. A DSL program embedding its values has no stable
  identifier, no template to cache, no separable sub-program to compile
  ahead of time. Separability is necessary, not optional.
  \emph{(Verification: logical / structural --- a script embedding its
  values produces a different cache key per invocation, making caching
  structurally pointless. The runtime documents itself: ``Un script
  funciona como F(x1,x2,\ldots,xn)'' (\texttt{ActorHandler.cs:1085}) ---
  F is the program; the xᵢ are its externalized values.)}
\item
  \textbf{DSL programs admit a structural binary taxonomy.}
  \emph{Parametric scripts} are separable from their values: cacheable
  as reusable programs with action identifiers, persistable as compact
  action entries in the journal, and economically compilable.
  \emph{Literal scripts} are non-separable: never cached, ephemeral, and
  persisted as their literal text. Whether either class is compiled or
  interpreted at runtime is governed by a separate per-actor policy
  that, by default, follows the parametric/literal split.
  \emph{(Verification: ActorHandler.cs:1091-1093 documents the rule
  verbatim. Implementation: PrepareCommandProgram
  (ActorHandler.cs:1097-1148) decides via parameters.HasUserParameter()
  and assigns one of three JournalEntry values --- IsScript (literal),
  IsNewAction (parametric, first invocation), IsExistingAction
  (parametric, cached invocation). Compilation flag set per program by
  AdjustCompilationMode (Program.cs:134-150) under the
  CompilationModePolicy enum.)}
\item
  \textbf{Interpretation and compilation are strategies over a shared
  AST substrate.} They are not separate execution paths over translated
  representations. Both walk the same AST: interpretation walks it
  directly; compilation specializes it into IL via Expression trees.
  Both paths pay the same parsing cost; only the compiled path pays the
  compilation cost. \emph{(Verification: Program.Execute()
  (interpretation) and Program.ExecuteExpression() →
  ProgramExpression().Compile() (compilation) operate over the same
  Statement tree produced by Parser.Parse().)}
\item
  \textbf{Journal density emerges from separability, not from a
  compression scheme.} A stable F(x₁, \ldots, xₙ) collapses repeated
  invocations into (actionId, values) tuples, with the script definition
  persisted once and referenced thereafter. The compactness is
  structural, not designed. \emph{(Verification:
  Diary.WriteNewActionEntry (definition + first invocation values) vs
  Diary.WriteActionEntry (actionId + values only on subsequent
  invocations).)}
\item
  \textbf{Verb richness: surface vs depth.} A single DSL verb may invoke
  orchestrations richer than its surface signature suggests. This is not
  the result of translation between layers; the DSL is the invocation
  language for domain operations implemented directly in the host
  language. Journal density is therefore the asymmetry between small
  persisted tokens and the live-domain operations they represent.
  \emph{(Verification: §4 develops the construct and §5.5 measures it on
  two open-source DDD aggregates --- eShop's \texttt{Order} purchase
  verb dispatches 9 host-language invocations with a 24-method static
  closure (β/α ≈ 2.7×); Grzybek's \texttt{SubscriptionPayment} verb
  dispatches 2 with a 73-method closure (β/α ≈ 36×). The magnitude of
  the asymmetry is a function of how much the host's persistence anchor
  compresses its domain, not of the runtime mechanism.)}
\item
  \textbf{Hot-loaded DSL programs over a stably-loaded domain.} The
  runtime supports hot-loaded DSL programs against a stably-loaded
  domain library: the assemblies configured as the actor's domain
  libraries are reflectively cached at first use; new DSL scripts can
  combine the types they carry in new ways without assembly reload. This
  enables ad-hoc procedure invocations against a live, stateful actor
  --- without redeploying endpoints or altering the domain library.
  \emph{(Verification: DomainLibraries.GetOrLoad(params Assembly{[}{]})
  caches public types statically per deduplicated assembly set;
  ActorV2.Using(scriptForChk, scriptForCmd) introduces new scripts per
  invocation against that cached surface area.)}
\end{enumerate}

\subsection{When NOT to use this
approach}\label{when-not-to-use-this-approach}

The principle of program--value separability has limits, and the
implementation pattern that realizes it has narrower limits still. We
name six regimes in which the analysis presented here either does not
apply, applies only partially, or invites overreach.

\subsubsection{A. This paper is not advocating for compiling all DSL
programs
unconditionally}\label{a.-this-paper-is-not-advocating-for-compiling-all-dsl-programs-unconditionally}

The structural taxonomy described here distinguishes parametric from
literal scripts; both are first-class. Literal scripts --- ad-hoc oracle
invocations against a live actor, exploratory queries against a stateful
runtime, configuration-time experiments --- pay zero amortization cost
because they incur no compilation. Treating them as legacy or transient
would discard a structurally valid mode of interaction.

\subsubsection{B. This paper does not address consistency under
concurrent
actors}\label{b.-this-paper-does-not-address-consistency-under-concurrent-actors}

The runtime described operates with strict per-actor ordering and
isolation. Multi-actor consistency, reaction propagation across actor
boundaries, and the contract under which observable state remains
coherent are treated formally in a companion paper of this series.

\subsubsection{C. This paper does not claim program--value separability
is sufficient for the four-fold
coherence}\label{c.-this-paper-does-not-claim-programvalue-separability-is-sufficient-for-the-four-fold-coherence}

Separability is necessary but not sufficient. SQL prepared statements
achieve compilation and caching while exhibiting only two of the four
faces --- the persisted artifact is row data, not the parameterized
statement itself, so homoiconic persistence does not emerge. Other
realizations may exhibit further subsets. Puppeteer demonstrates one
realization where all four faces emerge by structural commitment; the
converse is not entailed by separability alone.

\subsubsection{D. This paper does not extend separability to
general-purpose programming
languages}\label{d.-this-paper-does-not-extend-separability-to-general-purpose-programming-languages}

The principle is formulated for DSL runtimes --- domain-specific
languages whose programs are short, structured, and amenable to
identification by their textual form. General-purpose programming
languages introduce variables captured from enclosing scope,
side-effecting state in the host, and identity that exceeds the
parameter contract. Whether separability admits a meaningful formulation
in that broader setting is a separate question, and we do not take it up
here.

\subsubsection{E. This paper does not claim that hot-loaded DSL programs
replace traditional
deployment}\label{e.-this-paper-does-not-claim-that-hot-loaded-dsl-programs-replace-traditional-deployment}

Claim 6 articulates a runtime capability, not a deployment
recommendation. Hot-loaded DSL programs against a stably-loaded domain
enable ad-hoc invocation; redeployment of the domain library remains the
appropriate move when the implementation itself changes. The two are
complementary, not alternatives.

\subsubsection{F. This paper does not address security, sandboxing, or
resource bounds for hot-loaded
scripts}\label{f.-this-paper-does-not-address-security-sandboxing-or-resource-bounds-for-hot-loaded-scripts}

Claim 6 grants a running actor the ability to accept and execute DSL
programs formed at call time. The structural argument made here is
silent on the operational concerns that capability raises in production:
sandboxing of the DSL evaluator, per-script resource limits (CPU,
memory, wall-clock), authentication and authorization over which callers
may submit ad-hoc invocations, and the trust model under which the
supplied script is admitted at all. These concerns require separate
machinery --- out of scope for the conceptual contribution offered here,
but unavoidable for any deployment that exposes hot-loaded invocation
beyond a closed boundary.

\begin{center}\rule{0.5\linewidth}{0.5pt}\end{center}

\subsection{1. Introduction}\label{introduction}

This paper makes a design theory contribution. It identifies a
structural property of DSL programs that prior literature has touched
but not isolated as one --- the construct: \emph{program-value
separability}; derives the runtime consequences that follow when the
property is held --- the principles: compilation, caching, dense
journaling, replication-bounded entropy; and presents an instantiation
--- a system in which those principles have been realized in production
--- as confirmation that the construct is realizable. The contribution
is conceptual; the instantiation is the existence proof of
realizability, not the substance of the claim. The genre is the one
Hevner, March, Park, and Ram (2004) name design science research:
empirical evidence is presented in the form of a working artifact rather
than a controlled experiment.

\subsubsection{1.1 Calibrating the
surface}\label{calibrating-the-surface}

You already understand this taxonomy. A SQL prepared statement is
parameter-separable: the database parses it once, builds a query plan
once, and reuses both for thousands of invocations with different bind
values. An ad-hoc query offers no such surface --- each is a fresh
string, parsed and planned afresh. The same structural distinction
governs DSL programs in any runtime that aspires to compilation,
caching, and dense persistence. What changes between SQL and the runtime
described here is not the principle, but the surface to which it
applies.

That surface tends to be unfamiliar in a specific way. The professional
reader of this paper has very likely spent a career persisting
application state in relational tables. Where event sourcing has
appeared in their work, it has typically appeared in tabular form:
events as rows, schemas as projections, queries as joins against
materialized views. The implicit assumption --- that every
representation of state ultimately resolves into something table-shaped
--- is not a failure of imagination but the horizon that industrial
practice has produced over several decades. A runtime in which the
persisted artifact is the program itself, in which density emerges from
separability rather than from compression of rows, has had no canonical
place in that landscape.

Calibration begins by separating two layers that the dominant paradigm
tends to collapse into one. The first layer is the implementation:
domain types, methods, business rules --- written in the host language,
indistinguishable in form from any well-designed object-oriented or
functional library. The second layer is the invocation: the language by
which a running actor is commanded and queried. In runtimes of the kind
described here --- of which Puppeteer, the framework introduced in the
prior paper of this series, is one --- that invocation language is a
small DSL whose programs read as scripts, persist as journal entries,
and execute as either compiled lambdas or interpreted walks over their
own AST. The two layers are not translations of one another. The host
language carries the implementation; the DSL carries the contract by
which the implementation is reached. Conflating them --- assuming the
DSL must duplicate, transcribe, or shadow the host code --- is where
most readings of the architecture begin to diverge from its actual
structure.

Three readings should be pre-empted at the outset: this is not double
programming, not a translation layer, not a mapping specification. The
domain implementation lives once, in the host language. The DSL is the
language by which that implementation is invoked. The relationship is
closer to that between SQL and a relational engine than to that between
a model and a generated DTO: SQL does not duplicate the engine's logic;
it names operations that the engine performs. A short SQL statement can
trigger joins, locks, and execution plans far heavier than its surface
text. The DSL of an actor runtime carries the same kind of leverage ---
small surface, deep behavior --- and the persisted journal records not
the unfolding of that behavior but the invocations that produced it.
With the surface calibrated, the formal question can now be asked: what
does it take, structurally, for a DSL program to be compilable,
cacheable, and amenable to dense persistence?

\subsubsection{1.2 A formal
characterization}\label{a-formal-characterization}

The question can be sharpened. Compilation requires a stable artifact to
specialize. Caching requires a stable identifier under which to file the
result of that specialization. Dense journaling --- persistence that
does not balloon with the size of every invocation --- requires a stable
description of the operation that can be referenced rather than copied.
Each of these three demands the same property of the program: that it
admit, at the moment it is named, a separation between the program
itself and the values it will receive.

This property is \textbf{Program--Value Separability}. A DSL program is
separable when its textual form names an operation parameterized by
values supplied externally at invocation, rather than carrying those
values within its body. Separability is structural, not stylistic: it is
a property the program either has or does not have. It is decidable
syntactically at preparation time, by inspection of the parameter
declarations on the program's surface. The principle that follows admits
a compact statement: a necessary condition for cacheable specialization,
stable-key caching, and dense journaling of any DSL program is
program--value separability. Necessity is the strong claim; the limits
of what separability does not entail are taken up shortly.

Separability admits a binary structural taxonomy of DSL programs. A
\textbf{parametric script} is separable from its values: it is filed
under a stable identifier in a runtime cache and persisted to the
journal as a compact reference plus an argument vector. A
\textbf{literal script} is non-separable: it is not cached, runs once as
written, and persists as its literal text. Whether either class is
compiled or interpreted at runtime is governed by a separate per-actor
policy that, by default, follows the parametric/literal split --- but
admits override. Under that default policy, the runtime documents the
joint regime in its own comments: scripts without user parameters are
interpreted, uncached, and persisted as Script entries; scripts with
user parameters are compiled, cached with an ActionId, and persisted as
Action entries. The two strict properties --- caching and journal entry
format --- are encoded as values of a single enum decided at the moment
a program enters preparation. The taxonomy is not analytical but
operational: the runtime acts on it.

Necessity is not sufficiency. A program that admits separability gains
the structural possibility of being cached, persisted as a compact
reference, and amortizing its compilation --- but does not by that fact
alone exhibit dense journaling. SQL prepared statements illustrate the
gap: they are separable, they cache, they compile --- yet the persisted
artifact in the database is row data, not the parameterized statement
itself. The third face --- a journal in which the named operation is the
persisted entry, not its downstream effects --- does not emerge from
separability alone. It requires the further commitment, code-as-data,
that the runtime treat its programs as the persisted artifact: what is
recorded in the journal is the program parameterized and the values it
received, not the state changes those values produced. The sections that
follow trace how separability, joined to code-as-data, yields the three
faces in turn --- and where a runtime that holds separability without
the second commitment falls short.

The DSL program characterized above --- the textual artifact with
declared parameters --- is the unit of separability. When the term
\emph{program} appears at the substrate level in subsequent papers of
this series (the unit that is recorded, replicated between nodes, and
replayed against state), it denotes the pair: a \textbf{domain library}
(the compiled classes and verbs against which the DSL is interpreted,
loaded statically and identical across replicas) and a \textbf{journal
of invocations} (the ordered sequence of action references with
parameter bindings that the runtime has applied). Replay is the
deterministic re-execution of the journal against the library; what
persists, replicates, and reconstructs state is this pair. A DSL program
in the sense of §1.2 is the unit; the substrate-level program is the
cumulative state-producing artifact built from many DSL invocations
executed against a stable library. The two readings are not in tension:
separability of the unit is what makes the cumulative artifact dense;
the cumulative artifact is what makes the unit's separability
operationally consequential.

\subsection{2. Consequences of
separability}\label{consequences-of-separability}

\subsubsection{2.1 Compilation as
specialization}\label{compilation-as-specialization}

The first consequence of separability is that the program admits
ahead-of-time specialization. A separable program is, in form, a
function awaiting its arguments --- its body refers to values that will
arrive at invocation, not to literals fixed at write time. That
structural shape is precisely what a specializer requires. Given the
program once, the runtime can resolve its references, lower its abstract
syntax tree into a typed expression tree, and emit the executable form
in advance of any specific call. The values are bound later; the work of
preparing the program is done once.

The pipeline that produces this specialized form proceeds in three
stages. The parser yields an abstract syntax tree in which user
parameters appear not as values but as named slots --- placeholders to
be supplied at invocation. The runtime then traverses that tree and
emits a typed expression: each named slot is bound to a parameter
expression of the host runtime, and each operation of the script is
bound to the corresponding host-language method or property access.
Compilation of that expression yields an executable delegate over the
parameter list --- a function that takes the values supplied at
invocation and runs the program against them.

Specialization, considered alone, is a one-time cost paid against an
indefinite number of invocations. The parser runs once for the program;
the expression tree is constructed once; the compiler emits the delegate
once. Each subsequent invocation supplies new values into the existing
delegate and runs through host-language code that has already been
resolved, typed, and emitted. The amortization is structural: a
separable program admits compilation that pays for itself across
repeated invocation, while a non-separable program offers nothing to
amortize against. The work, however, is only economical if the
specialized form persists between invocations.

\subsubsection{2.2 Caching as
identification}\label{caching-as-identification}

Persistence between invocations requires an identifier under which to
file the work that has been done. Separability supplies that identifier
directly: the script string of a separable program is stable across
invocations because its values are external to it. The same
parameterized program text, encountered a second time, is recognizable
as the same program --- and the cached specialization belongs to it. A
non-separable program has no such anchor; every invocation produces a
different string, and no two invocations are the same program at all.
Caching is not an addition; it is what becomes possible once the program
has a stable name.

The mechanism is direct. On first encounter, the runtime parses the
script, prepares its specialized form, and stores both under an action
identifier --- a stable handle assigned to the script when it is first
encountered. The script string serves as the lookup key; the action
identifier serves as the persistent reference under which the program's
specialization lives. On subsequent encounters with the same script, the
lookup succeeds, and the runtime retrieves the specialization without
repeating the work that produced it. The cache is not a memoization of
values; it is a registry of programs.

On a cache hit, the work that has been preserved is the program itself:
the resolved tree, the typed expression, the compiled delegate. The
values arriving at this invocation are bound into the delegate's
parameter slots and the program runs against them. From the program's
standpoint, nothing has changed between the first invocation and the
thousandth --- the script names the same operation, parameterized by the
same slots, executed against the same domain. What differs across
invocations is only the values supplied; what persists is the program.

\subsubsection{2.3 Dense journaling as
reference}\label{dense-journaling-as-reference}

The third consequence of separability is that persistence becomes
reference rather than duplication. The program, having earned a stable
identifier in the cache, can be written to the journal once --- as a
definition, with the script's text and its parameter declarations
recorded in full. Each subsequent invocation is recorded not by
duplicating the program text but by referring back to that definition
under its identifier, accompanied only by the values supplied at that
invocation. Where a non-separable runtime must record every invocation
as a complete script --- body, values, and all --- a separable runtime
records the program once and its invocations as compact references
against it.

The mechanism distinguishes three kinds of journal entries. When a
parametric program is encountered for the first time, the runtime writes
a definition entry: the script's text, its parameter declarations, and
the values supplied at this first invocation, all recorded together
under the program's newly assigned identifier. Subsequent invocations of
the same program write only a reference entry --- the identifier of the
definition, plus the values for that call. A non-separable program,
having no stable identifier, is written each time as a script entry: its
text in full, since it cannot be referred back to anything that came
before. Three entry kinds, decided mechanically at preparation time from
the program's separability, encode the program's relation to its
identity.

The density follows. For a parametric program invoked many times, the
journal's storage scales with the cumulative size of the argument
vectors, not with the cumulative size of the script's body --- the body
has been recorded once, and every later invocation costs only the bytes
of its values. This is what makes the journal dense: nothing is copied
that has been preserved by reference, and what survives is the operation
that took place, not the state that resulted. The persisted artifact is
the program --- the script as text, parameterized, with its invocation
values --- rather than the downstream effects of running it. In
code-as-data terms, the program lives in the journal as itself, not as a
record of what it produced. Compilation, caching, and dense journaling
are thus not independent optimizations but successive consequences of
the same structural property: once a program is separable from its
values, it becomes nameable; once nameable, it becomes cacheable; once
cacheable, it becomes referable in persistence. The mechanism by which
the runtime accomplishes this --- the pipeline that compiles, the cache
that retains, the journal that references --- is the subject of §3.

\subsection{3. Realization in a concrete
runtime}\label{realization-in-a-concrete-runtime}

\subsubsection{3.0 Origins of the
instantiation}\label{origins-of-the-instantiation}

The realization described in this section was not engineered to
demonstrate the principle articulated above. The principle was
identified by inspecting a runtime that had, over years of independent
development for production use, come to satisfy it. The genealogy is
structural rather than programmatic --- first the artifact, then the
principle that explains why the artifact cohered. What follows describes
the runtime as it stands; the relationship between its mechanisms and
the principle of §1.2 is a recognition, not a derivation.

\subsubsection{3.1 Compilation pipeline}\label{compilation-pipeline}

The mechanisms described here are not independent engineering decisions
but direct realizations of the separability principle established in the
preceding sections. In the runtime described, the compilation pipeline
runs from text to executable in three well-defined stages. A
\texttt{Parser.Parse()} call lowers the script into an abstract syntax
tree of \texttt{Statement} and \texttt{AstExpression} nodes. Each
\texttt{Statement} carries an \texttt{ExecuteExpression} method that,
when traversed, emits a \texttt{System.Linq.Expressions} node bound to
the host parameters and the host-language operations the script names.
The runtime composes these into a single \texttt{Expression.Lambda} and
calls \texttt{.Compile()} to produce an executable delegate. Three
stages --- parsing, lowering, and compilation --- produce the artifacts
on which the runtime operates: the AST, the expression tree, and the
compiled delegate.

The mechanism is straightforward to follow in the source.
\texttt{PrepareCommandProgram} (\texttt{ActorHandler.cs:1097-1148})
orchestrates the first encounter: it rents a parser from a pool, parses
the script, and decides --- by parameter presence --- whether to enter
the parametric path. On the parametric path, the program's
\texttt{ProgramExpression} method (\texttt{Program.cs:182-244}) walks
the AST: for each Statement, it calls the Statement's
\texttt{ExecuteExpression}, which returns an \texttt{Expression} node
bound to the runtime's parameter and output expressions. These nodes
accumulate into a \texttt{BlockExpression}, which
\texttt{ProgramExpression} wraps in
\texttt{Expression.Lambda\textless{}Func\textless{}Parameters,\ Output,\ string\textgreater{}\textgreater{}}
(\texttt{Program.cs:244}). The compilation itself is one line:
\texttt{\_executable\ =\ programExpression.Compile()}
(\texttt{Program.cs:163}). The compiled delegate lives in the program's
\texttt{\_executable} field; subsequent invocations skip the compile
step and call the delegate directly with the values supplied at the call
site.

The same mechanism applies one level deeper, to a feature the runtime
calls \emph{Eval parameters}. A parameter declared with the
\texttt{Eval} modifier carries its own sub-script, which the runtime
treats as a separable program in miniature: parsed once, compiled to its
own delegate, and cached against the parameter's name. The cache
structure is a dictionary keyed by parameter name to a tuple of the eval
script and its compiled executable (\texttt{Program.cs:305}). On each
invocation, the runtime compares the parameter's current script to the
cached one; if they differ, the entry is invalidated and the sub-program
is re-parsed and re-compiled (\texttt{Program.cs:323-331}). The
principle scales: any region of the program that has a stable textual
identity admits the same treatment as the whole.

\subsubsection{3.2 Interpretation
retained}\label{interpretation-retained}

The runtime preserves interpretation as a first-class execution mode,
not as a legacy fallback. When a script presents itself without user
parameters, or when a per-actor policy directs the runtime to interpret
regardless of parameter presence, the program executes by walking the
AST directly: each Statement carries an \texttt{Execute} method that
runs against the program's current state, with no expression tree built
and no delegate compiled. The interpreted path costs zero compilation
but pays the cost of tree traversal at every invocation. Its place in
the runtime is precisely the inverse of the compiled path's: useful
where invocation is one-shot or unanticipated, useless where the program
will be reused.

The decision is governed by a per-actor \texttt{CompilationModePolicy}
enum with three values --- \texttt{Automatic}, \texttt{AlwaysCompiled},
and \texttt{AlwaysInterpreted} --- declared on the actor and defaulted
to \texttt{Automatic} (\texttt{Actor.cs:8-13}). On the parametric path,
\texttt{PrepareCommandProgram} calls \texttt{AdjustCompilationMode} with
\texttt{useInterpretedMode:\ false}; on the literal path, with
\texttt{useInterpretedMode:\ true} (\texttt{ActorHandler.cs:1117-1131}).
Under \texttt{Automatic}, the policy honors that hint: parametric
programs compile, literal ones interpret. Under \texttt{AlwaysCompiled}
or \texttt{AlwaysInterpreted}, the hint is overridden in favor of the
policy's name (\texttt{Program.cs:134-150}). Execution itself is
dispatched by another switch on the same policy: \texttt{Perform} calls
\texttt{ExecuteExpression} if the program is in compiled mode and
\texttt{Execute} otherwise (\texttt{ActorHandler.cs:1955-1968}).

Caching applies asymmetrically to the two principal kinds of DSL
invocation. For commands --- programs that mutate actor state and
persist to the journal --- the runtime caches only when the script
declares user parameters (\texttt{ActorHandler.cs:1117}). For queries,
checks, and emit invocations --- programs that read state without
persistence --- the runtime caches when the script declares user
parameters or when no parameters are supplied at all
(\texttt{ActorHandler.cs:1405}). The asymmetry is operational. Commands
run under a write lock and serialize, so caching pays only when the same
parametric program is invoked many times. Queries run under a read lock
and parallelize, so a cache entry amortizes regardless of parametric
reuse --- a frequently-emitted query gains from caching even when its
parameter set is fixed. One final observation: cosmetic variation in the
script's text across releases generates distinct cache identifiers, but
the runtime's reaction mechanism --- treated in a companion paper ---
matches behavior against semantic patterns rather than identifier
equality, so coherence is preserved across textual variants.

\subsubsection{3.3 Hot-loaded DSL programs over a stably-loaded
domain}\label{hot-loaded-dsl-programs-over-a-stably-loaded-domain}

The third element of realization is the runtime's support for hot-loaded
DSL programs against a stably-loaded domain library. The phrase deserves
care, because it is easy to misread. The hot element is the DSL: at any
time during the actor's life, a new script --- never seen before by this
runtime instance --- can be supplied, parsed, prepared, and executed
without restart, redeployment, or reflection over a re-loaded assembly.
The stable element is the domain library: the host-language types and
methods that the DSL invokes are loaded once, by reflection over the
assemblies configured as the actor's domain libraries, and held in a
cache for the lifetime of the process. Hot-loaded programs combine
stable domain elements in new ways; they do not bring new domain
elements into being. The runtime is hot at one layer and stable at the
layer beneath it.

The mechanism is brief in the source. When an actor is constructed, the
public types of its configured library assemblies are walked once by
reflection and cached against a deduplicated key:
\texttt{DomainLibraries.GetOrLoad(params\ Assembly{[}{]})}
(\texttt{DomainLibraries.cs:77-115}, with both a single-assembly and a
multi-assembly overload), invoked from \texttt{ActorHandler}
(\texttt{ActorHandler.cs:61}) with the actor's
\texttt{LibraryAssemblies}. The cache survives the lifetime of the
process; the same assembly set, loaded by a second actor, reuses the
same domain dictionary. Against that fixed surface, new DSL scripts
arrive through the actor's fluent interface ---
\texttt{ActorV2.Using(scriptForChk,\ scriptForCmd)}
(\texttt{ActorV2.cs:32}) --- at any time the actor is running. There is
no AppDomain reload, no \texttt{AssemblyLoadContext} unload, no
file-watcher over the assembly: the surface area is fixed at first load
and dynamics live entirely above it.

The operational consequence is direct: a running actor can be queried or
commanded with a script formed at call time, against any combination of
the domain operations the configured library assemblies carry. A
configuration adjustment, a one-off audit query, a derivation that the
published endpoints do not anticipate --- each can be expressed as a
fresh DSL program and executed against live actor state without
intermediate deployment. The relationship resembles gRPC in its
directness --- a procedure call against the live system rather than a
navigation of REST resources --- but without a pre-declared interface
contract: the script itself is the contract, formed at call time. What
has been shown in this section is that separability is not merely a
formal property of DSL programs but one that a concrete runtime can
recognize, act upon, and encode into its execution, caching, and
persistence behavior. The realization characterized here is one of an
extensible family: other actor runtimes --- Akka, Orleans, Erlang/OTP
--- admit, in principle, the same construction, since a parameterized
DSL surface above a host-language domain library would yield the same
four-fold coherence under the same separability commitment. The
expressive reach of such on-the-fly programs depends on the richness of
the verbs the domain library exposes, which the next section takes up.

\subsection{4. Verb richness: surface vs
depth}\label{verb-richness-surface-vs-depth}

Separability makes the program nameable; verb richness determines how
much domain meaning that name carries. The verbs that a DSL invokes are
not method calls on data structures. A verb in this setting names an
operation against a live, stateful actor --- an orchestration that the
host language has been written to perform on behalf of the domain. A
single verb may invoke many internal methods, traverse domain
relationships, evaluate derived properties, trigger downstream
behaviors, and write to multiple regions of state, all in the course of
executing one DSL statement. Where a setter assigns a field, an actor
verb concludes a transaction --- moving the actor from one consistent
state to another. The asymmetry between the surface --- the few
characters that name the verb in the script --- and the depth --- the
volume of domain computation that runs when it is called --- is
structural, not incidental.

Consider a verb that any reader from an e-commerce or supply-chain
background will recognize: the \texttt{Confirm} of a purchase order. The
script that names this verb at the DSL surface is short --- only a
handful of tokens. What runs when the script is invoked is substantially
larger: a verification that the order's reservations are still valid, a
hold-to-allocation conversion across one or more inventory ledgers, a
tax computation that depends on the order's line items and their fiscal
contexts, a posting to the financial ledger, and a set of reaction
triggers that propagate the confirmation to subscribed views and
downstream actors. Each of these steps is itself a tree of host-language
method calls; the leaf operations include arithmetic, predicate checks,
persistence writes, and outbound messages. The runtime gives the
consistency boundary that separates this depth from its surface
first-class support: a check script and a command script can be supplied
together via \texttt{ActorV2.Using(scriptForChk,\ scriptForCmd)} and
dispatched as a single invocation through
\texttt{PerformCheckThenCommand} (\texttt{ActorV2Invocation.cs:69}). The
check runs against the actor's current state under a read lock; if its
assertions hold, the command runs under a write lock and is persisted to
the journal. If the check fails, no state change is applied and no
journal entry is written --- the actor's consistent state is preserved
by construction.

Confirm is one example among many. The same asymmetry obtains for any
verb that names a domain transition rather than a field assignment ---
Settle, Allocate, Reconcile, Release, and a hundred others that a
well-designed domain library will expose. Verb richness is not an
artifact of this particular runtime; it is a property of how domain
operations are conceived and named. The implication folds back into the
journal density argument of §2.3: when each persisted entry refers to a
domain operation that may run dozens to hundreds of internal method
calls, the journal's compactness in bytes is matched by an enormous
compactness in semantic information per byte. The journal does not
record state changes; it records named transactions, each of which
carries its full operational meaning by reference to the domain library
that knows how to run it. The value of separability would be modest if
the verbs it named were shallow. It becomes profound when each name
stands for a full domain transaction.

Under porous substrates, the asymmetry inverts entirely. There, every
domain element must persist into a relational schema, and deep domain
logic becomes a serialization burden --- each derived state, each
accumulator update, each new field propagates as schema complexity, ETL
transforms, and projection overhead. Under a code-as-data substrate, the
same depth is rewarded, not penalized: the domain library can be made
arbitrarily expressive without expanding what the journal must record.
Complex abstractions compose without friction: no DTOs, ORM annotations,
or projection layers stand between the domain operation and its
persistence. Porosity makes deep domains expensive to represent;
anti-porosity makes them economical to compose. The empirical magnitude
of these numbers --- how many operations a typical verb dispatches, how
many bytes a typical journal entry occupies, how the ratio behaves at
scale --- is the subject of §5.

\subsection{5. Empirical results}\label{empirical-results}

\subsubsection{5.1 Methodology}\label{methodology}

The measurements that follow span two complementary axes. The first axis
is \emph{program complexity}: a synthetic straight-line arithmetic
kernel parametric on integer inputs (depth 5 to 100 statements); a
DSL-rich kernel exercising control flow (for, if), arithmetic, and
parameter binding (\textasciitilde500 dispatched operations per
invocation); and a multi-item purchase command operating against an
external rich-domain aggregate. The first two tiers isolate runtime
overhead independent of host-language work; the third locates that
overhead within a representative end-to-end transaction.

The second axis is \emph{host codebase}. The production-verb
measurements are replicated against two independent open-source
MIT-licensed DDD aggregates that share a structural shape (rich
behavioral methods on aggregate roots) but represent disjoint business
domains: \texttt{dotnet/eShop}'s \texttt{Order} aggregate (e-commerce
ordering with multi-item cart and a four-step state machine), and
\texttt{kgrzybek/modular-monolith-with-ddd}'s
\texttt{SubscriptionPayment} aggregate (subscription billing with a
payment-lifecycle verb). The dual-codebase design is deliberate: if the
measured properties of compilation, caching, and journaling are
structural consequences of the runtime rather than artifacts of a
particular domain, the same property should appear on both codebases.
The selection criterion was structural similarity (DDD aggregates with
rich behavior methods rather than CRUD entities) over an unrelated
business domain. Both aggregates are reachable as public source for
reproduction; Appendix A lists the per-section datasets, the runtime
branch they were produced from, and the commit SHA.

All measurements use Stopwatch instrumentation around the runtime's
compilation and execution paths, with N≥1,000 invocations per cell after
warm-up runs are discarded; cache and pool hit rates use Interlocked
counters; journal-density measurements parse the runtime's binary
journal format directly. The bench follows a bootstrap-then-measurement
pattern: a non-parametric bootstrap script establishes the facade in the
actor's persistent symbol table once, and subsequent timed iterations
invoke a parametric measurement script that binds per-iteration values
via the runtime's parameter API. The script string stays constant across
iterations, so the compiled-program cache hits on every call --- the
regime the amortization argument names.

\subsubsection{5.2 Compilation
amortization}\label{compilation-amortization}

Compiled execution shows a p50 speedup over interpreted execution that
varies with where the program's work resides. For a synthetic arithmetic
kernel of depth 100, the speedup is 4.10× (N=1,000 invocations per
cell). For a DSL-rich kernel of \textasciitilde500 dispatched operations
exercising for-loops, conditionals, and parameter binding, the speedup
is 2.92×. For a multi-item purchase verb against the eShop
\texttt{Order} aggregate --- one DSL dispatch cascading through
\texttt{Order} ctor, four \texttt{AddOrderItem} calls, and a four-step
state-machine walk to Shipped --- the speedup is 1.49× (interpreted 5.5
µs p50, compiled 3.7 µs p50, N=1,000). The pattern is monotonic: the
more of the per-invocation work is DSL-bound, the larger the speedup;
the more it is domain-bound, the smaller. The runtime amortizes only the
AST traversal overhead --- the host-language code dispatched by the verb
runs identically in both modes.

Specialization is paid for at first encounter, not at every invocation.
For straight-line arithmetic kernels, the cost of compiling a single
program ranges from 0.8 ms (p50, scripts of 5-29 statements) to 4.1 ms
(p50, scripts of 80-104 statements); for richer DSL kernels exercising
control flow, from 1.1 ms to 5.3 ms over the same depth range. The cost
grows linearly in statement count --- approximately 44 µs per statement
for arithmetic, 43 µs per syntactic structure for richer kernels. For
the eShop purchase verb, the cold compile cost is 281 µs at p50 (394 µs
at p95, N=100 distinct script variants forcing cache misses). Against
the 1.8 µs per-invocation speedup over interpreted execution for the
same verb, this cost recovers itself after approximately 156
invocations. Any actor whose lifetime exceeds the first second of
operation amortizes its compilation investment in full. These numbers
are not performance claims but empirical confirmation of the structural
amortization predicted by separability.

\subsubsection{5.3 Cache amortization}\label{cache-amortization}

The same amortization pattern applies recursively to sub-programs. A
parameter declared with the Eval modifier carries its own sub-script
that the runtime treats as a separable program in miniature; on a cache
hit, the sub-program pays only the cost of invoking its cached delegate,
while on a cache miss --- when the sub-script's text differs from the
cached entry --- it re-parses, rebuilds, and re-compiles. Across three
sub-program complexities (a two-term arithmetic expression, a fifty-term
arithmetic kernel, and a method call against a facade-bound counter on
the eShop side), the cache-hit cost ranges from 1.7 to 2.2 µs at p50,
while the cache-miss cost ranges from 223 to 263 µs at p50. The ratio
between cache miss and cache hit sits at approximately 120× regardless
of sub-program complexity --- a two-orders-of-magnitude separation that
confirms the principle extends to any region of the program with a
stable textual identity.

Allocation pressure is amortized at a finer granularity by parser and
parameter pools. Three workloads measured under AlwaysCompiled policy:
1,000 single-thread invocations against a stable parametric script
recorded a 100\% parameter-pool hit rate; 1,000 single-thread
invocations against distinct cold-cache scripts recorded a 100\% hit
rate on both pools; 5,000 invocations across 8 parallel threads against
distinct scripts recorded 99.90\% on parsers and 99.86\% on parameters
--- twelve total misses across the parallel workload's
\textasciitilde10,000 pool rents. Allocation of new Parser and
Parameters instances is incurred at startup and under brief
thread-fanout transients only; under steady state, both pools service
the rent without allocation.

\subsubsection{5.4 Journal density}\label{journal-density}

In a parametric purchase workload against the eShop \texttt{Order}
aggregate --- a nine-line DSL script cascading through order
construction, four \texttt{AddOrderItem} calls, and a four-step
state-machine walk --- 99.8\% of journal entries land as compact action
references: 1,001 invocations produce 1,001 action entries plus a single
Define entry that carries the script body once. Each action entry's
payload averages 115 bytes --- the argument vector for seventeen
parameters (user identity plus four products' details plus shared
modifiers). The Define entry's payload is 642 bytes --- the script body
itself, persisted once. Had each invocation stored the literal script
text instead, the action payload would be 5.6× larger. The runtime's
choice to persist named operations rather than their textual content is
what makes this density possible. This density does not arise from
compression techniques but from reference instead of duplication.

The same compaction structure appears on the second host. Against
Grzybek's \texttt{SubscriptionPayment} --- a two-statement DSL surface
(\texttt{NewWaitingPayment} followed by \texttt{MarkAsPaid}) over a
payment-lifecycle aggregate --- N=1,000 parametric invocations produce
1,001 Action entries (99.8\% of the journal) plus a single Define entry
of 207 bytes carrying the script body. Action payloads average 60 bytes.
The literal-script storage would be 3.5× larger. The ratio is more
modest than eShop's because both the script body and the argument vector
are smaller on a narrower verb; the structural property --- arguments
scaling with invocations while the body is persisted once --- is
identical.

The magnitude of the ratio is a function of how the host's domain API
surfaces its data, not of the compaction mechanism. eShop's
\texttt{AddOrderItem} carries the full product specification per item
--- pid, name, price, discount, picUrl, units --- so per-iteration
argument vectors are dense (115 B). Grzybek's \texttt{NewWaitingPayment}
accepts five primitive parameters and \texttt{MarkAsPaid} takes none, so
the argument vector is short (60 B). Hosts whose verbs identify catalog
entries by short stable references compound the ratio further; hosts
whose verbs string together longer parametric sequences also compound
it. Either way the structural property holds: arguments scale with
invocations; the script body does not.

\subsubsection{5.5 Verb richness}\label{verb-richness}

For the eShop purchase verb the DSL script dispatches 9 host-language
invocations per invocation, exactly reproducible across 1,000 measured
runs. Static forward closure of the call graph from those 9 entry points
reaches 24 methods within \texttt{dotnet-eShop/src/Ordering.Domain} ---
measured by a syntactic walker that traces method-to-method edges from
the entry points, excludes trivial accessors, and treats the project
boundary as a structural ceiling. For the Grzybek
\texttt{SubscriptionPayment} verb the dispatch count is 2 (the facade
absorbs the value-object assembly into one host-language call from the
DSL's perspective) and the closure reaches 73 methods within the
project's \texttt{Payments.Domain} plus its shared
\texttt{BuildingBlocks/Domain} (275 declarations indexed, 105 trivial
accessors filtered). The β/α asymmetries are 2.7× for eShop and
\textasciitilde36× for Grzybek.

These two numbers do not represent the ceiling of the mechanism --- they
represent its floor. Both \texttt{Ordering.Domain} and
\texttt{Payments.Domain} are RDBMS-anchored DDD aggregates: their
authors designed against the anticipation of EF Core hydration, and that
anticipation flattens polymorphism, caps graph depth, and externalizes
references as foreign keys. The marks are directly visible in the source
--- \texttt{private\ set} on every property, explicit
\texttt{EF\ Core\ 2.0\ owned\ entity} annotations on value objects,
\texttt{int?} foreign keys instead of object references, and a
deliberate absence of polymorphism on the aggregate roots. Domains
developed without this pressure on their representation grow
polymorphism, deeper graphs, and richer cross-class behavior; the same
mechanism then traverses substantially further. The β/α magnitude is
therefore a measurement of how much the host's persistence anchor
compresses the domain --- not of the mechanism the runtime applies once
a verb is invoked. The mechanism is what §4 anticipates: a verb whose
surface signature names an orchestration deeper than its syntax
suggests. The depth of that orchestration is a property of the host's
domain, not of the runtime that dispatches into it.

\subsection{6. Counter-arguments}\label{counter-arguments}

\subsubsection{\texorpdfstring{6.1 \emph{``This is just JIT
compilation''}}{6.1 ``This is just JIT compilation''}}\label{this-is-just-jit-compilation}

A reviewer familiar with managed runtimes might object that the runtime
described here is just-in-time compilation in disguise. The objection
conflates two distinct compilation events. The .NET CLR's JIT translates
IL bytecode into native machine code at first invocation of any method;
this happens identically whether the runtime executes its DSL by walking
an AST or by invoking a compiled delegate. What the runtime adds is a
separate, prior compilation: from DSL script to typed Expression tree to
IL --- the work that produces the delegate the JIT will eventually
translate. The speedup measured in §5.2 is a function of this DSL→IL
stage, not of the IL→native stage. Both modes deliver IL to the JIT in
the same way, so JIT effects cancel; what remains is the AST traversal
overhead that the DSL→IL stage amortizes away. Calling this ``just JIT''
would erase the layer of specialization that separability makes
possible.

\subsubsection{\texorpdfstring{6.2 \emph{``Why retain interpretation?
Just compile
everything''}}{6.2 ``Why retain interpretation? Just compile everything''}}\label{why-retain-interpretation-just-compile-everything}

A second objection: if compilation is so much faster, why retain
interpretation at all? The answer is that compilation is economical only
when the program will be reused. A script seen once and never again ---
an ad-hoc query against a live actor, a one-off configuration
adjustment, an exploratory invocation formed at call time --- provides
no future invocations against which to amortize the compilation cost.
Forcing such scripts through the compilation pipeline pays the 1.4 ms
cold cost (§5.2) and discards the result; the compiled delegate is never
invoked a second time, and the journal records the script literally
rather than as a referenced action. The runtime's
\texttt{AlwaysCompiled} policy permits this on demand, but the default
policy correctly recognizes that not every script is amortizable.
Separability is necessary for compilation to make sense; reuse is the
condition under which compilation is economical.

\subsubsection{\texorpdfstring{6.3 \emph{``Dual paths add memory
cost''}}{6.3 ``Dual paths add memory cost''}}\label{dual-paths-add-memory-cost}

A third objection: maintaining a compiled delegate per parametric
program inflates memory at scale, and the dual-path arrangement
compounds the cost. The measurement disagrees. In a single-actor cache
holding 100, 1,000, 10,000, and 100,000 distinct parametric programs
under \texttt{AlwaysCompiled} policy, the per-entry footprint stabilizes
at approximately 6 KB across all four scales: 6,039 bytes per entry at
100 programs, 5,981 at 1,000, 5,977 at 10,000, and 5,930 at 100,000 ---
no super-linear overhead, no hash-table degradation. The marginal memory
of an invocation against an already-cached program is effectively zero:
100,000 invocations against a single cached program retain 104 bytes
total. A production actor caching dozens of distinct parametric programs
and invoking them at scale incurs a memory cost on the order of hundreds
of kilobytes. The cache scales linearly with distinct programs, not with
invocations --- well below the threshold at which the dual-path
arrangement would be a structural concern.

\subsubsection{\texorpdfstring{6.4 \emph{``This is just SQL prepared
statements''}}{6.4 ``This is just SQL prepared statements''}}\label{this-is-just-sql-prepared-statements}

A final objection: this is just SQL prepared statements rebranded ---
the same parametric-template-plus-bind-values pattern that relational
databases have used for decades. The pattern is indeed the same; what
differs is what the runtime persists. A SQL prepared statement is
parameter-separable and admits compilation and caching: the database
parses it once, builds a query plan once, and reuses both for many
invocations with different bind values. But what the database persists
is row data --- the projections produced by executing the statement
against the underlying tables. The parameterized statement itself is not
the persisted artifact; the rows it touches are. The runtime described
here persists the operation rather than its row-level effects. Two of
the four faces of separability --- compilation and caching --- are
present in both systems; the journal density that follows from
persisting the named operation is what SQL prepared statements do not
achieve. Separability is necessary; pairing it with code-as-data is what
produces the dense journal.

\subsubsection{\texorpdfstring{6.5 \emph{``A verb dispatching dozens or
hundreds of host methods produces an opaque
runtime''}}{6.5 ``A verb dispatching dozens or hundreds of host methods produces an opaque runtime''}}\label{a-verb-dispatching-dozens-or-hundreds-of-host-methods-produces-an-opaque-runtime}

A reviewer might object that a verb whose static call-graph closure
reaches dozens or hundreds of host-language methods (§5.5) produces an
opaque runtime, and that debugging across such depth is intractable. The
objection misreads the architecture. The depth lives in the host
language: domain classes, methods, business rules --- debuggable with
the standard tools the host already provides. The DSL surface only names
what the host invokes; it does not introduce a separate execution layer
that the host debugger fails to penetrate. Standard practice in the
runtime described here proceeds by test-driven development with
end-to-end test cases, with breakpoints and step-through performed in
the host language; once a script is moved to a runtime endpoint, the
same breakpoints continue to apply against the domain library it
invokes. More consequentially, the journal makes post-mortem debugging
tractable in a way porous substrates do not allow. When a production
failure surfaces at journal entry id N, that entry refers --- by name
--- to the exact script that produced the defect, parameterized by the
exact values the script consumed. A breakpoint at the journal entry, a
step-through of the named operation, and the defect reproduces; the
failure is then replicated in a small end-to-end test case, fixed, and
released. The dense journal is not the opacity it might appear to be on
a superficial reading; it is among the runtime's most powerful debugging
artifacts, because the persisted artifact is the operation that ran, not
a downstream trace of its effects.

Each objection treats compilation, caching, and journaling as
independent engineering choices. The argument of this paper is that they
are consequences of a single structural property; the objections
dissolve when that property is made explicit.

\subsection{7. Related work}\label{related-work}

\subsubsection{7.1 Partial evaluation}\label{partial-evaluation}

The structural condition this paper formalizes has a close ancestor in
\textbf{partial evaluation}, as developed by Jones, Gomard, and Sestoft
(1993) and the Futamura projections (Futamura 1971, 1999). Partial
evaluation specializes a program with respect to a subset of inputs
known in advance, producing a residual program that depends only on the
dynamic arguments --- and it requires, structurally, that the program's
static and dynamic inputs be separable. Where the partial-evaluation
literature applies the technique to general-purpose host languages, this
paper applies the same separability principle to a domain-specific
language whose programs are short, parametric by construction, and
stored as journal entries; in this setting, separability is syntactic,
decidable from the program's surface rather than from static analysis.
The further consequences traced in §2 --- caching and dense journaling
--- are not part of the partial-evaluation tradition; they follow from
pairing separability with code-as-data persistence.

\subsubsection{7.2 Lambda lifting}\label{lambda-lifting}

\textbf{Lambda lifting} (Johnsson 1985) is a program transformation that
rewrites locally-defined functions whose free variables capture an
enclosing scope into top-level functions that receive those variables as
explicit parameters. The transformation makes the function's
dependencies syntactically visible --- and, by doing so, makes the
function amenable to ahead-of-time compilation, separate compilation,
and uncluttered call-graph analysis. The structural insight is the one
this paper generalizes: dependencies that remain implicit in the
program's body cannot be specialized over, but the same dependencies
promoted to the surface can. Where lambda lifting transforms programs
that were already written and externalizes their captured environment,
the runtime described here requires programs to declare their values as
externalized parameters from the start. Lambda lifting opens
compilation; this paper extends the same principle to caching, dense
journaling, and replication-bounded entropy by joining separability with
code-as-data persistence.

\subsubsection{7.3 Closure conversion}\label{closure-conversion}

\textbf{Closure conversion} (Appel 1992) is a compilation technique that
translates a function with captured lexical environment into an explicit
closure record --- a pair of function pointer and environment data ---
that the runtime carries as a single value. Closure conversion makes
implicit lexical dependencies explicit, but preserves them as runtime
data: the closure travels with its environment record, retained in
memory, and re-bound on each invocation. Lambda lifting and the
principle of this paper sit on the other side of the same choice: rather
than carry the environment, eliminate it by promoting the variables that
would have been captured into top-level parameters supplied at the call
site. The runtime described here makes that choice mandatory and
declarative --- values must be supplied at invocation, not captured from
a surrounding scope --- which is what permits the journal to record the
program by name without any captured-environment record traveling
alongside it.

\subsubsection{7.4 Template instantiation}\label{template-instantiation}

\textbf{Template instantiation} in general-purpose languages ---
exemplified by C++ templates (Stroustrup 1986--present; Vandevoorde and
Josuttis 2017) and the broader tradition of generic programming ---
applies the separability principle in another setting: a template body
names operations parameterized by types or compile-time constants, and
the compiler instantiates each template with concrete parameters at
compile time, producing specialized code per instantiation. The
structural pattern is the same as the one described here: a program that
admits parametric reuse must be written so its parameters are separable
from its body. Template instantiation operates over types and
compile-time values; the runtime described here applies the same
principle over runtime values supplied at invocation. The distinction is
operational: instead of monomorphizing the program once per parameter
tuple, the runtime compiles the program once and rebinds its arguments
on every call. Both arrangements presuppose what this paper formalizes.

\subsubsection{7.5 Prepared statements}\label{prepared-statements}

\textbf{Prepared statements} in relational database systems (SQL
standard ISO/IEC 9075; Hellerstein and Stonebraker 2007) are the closest
operational analog to the runtime described here. A prepared statement
carries a parametric template (with placeholders for bind values), the
database parses and plans it once, caches the plan, and reuses both for
many invocations with different argument vectors --- the same
parametric-template-plus-bind-values pattern §1.1 used as pedagogical
anchor for the principle this paper formalizes. The two systems share
separability and its first two consequences: compilation (the query
plan) and caching (the plan cache). They diverge on what the runtime
persists. A relational database persists row data; the prepared
statement itself is not the persisted artifact. The runtime described
here persists the named operation rather than its row-level effects,
completing the four-fold coherence with journal density and
replication-bounded entropy. The principle this paper formalizes
generalizes beyond any one runtime --- it characterizes the structural
condition any DSL runtime must satisfy to admit compilation, caching,
and dense persistence at all.

These traditions arise in different domains --- compilers, functional
languages, generic programming, database engines --- but they converge
on the same structural observation: programs whose dependencies are
syntactically separable from their bodies admit forms of specialization,
reuse, and compact representation that programs with embedded values
cannot. This paper isolates that observation as a single principle and
traces its full runtime consequences in a DSL setting.

\subsection{8. Conclusion}\label{conclusion}

The argument of this paper can be stated compactly. Program--value
separability --- the syntactically decidable property that a DSL program
declares user parameters rather than embedding values in its body --- is
the structural precondition under which compilation, caching, and dense
journaling become economically meaningful. The four runtime faces traced
here are not parallel design choices: they are downstream consequences
of separability, paired with code-as-data persistence in the case of
dense journaling. The runtime characterized in §3 realizes the principle
concretely; the magnitudes reported in §5 confirm the predicted
structural amortization across two open-source DDD aggregates --- a
compiled-versus-interpreted speedup that scales monotonically with
DSL-bound work (1.49× on the eShop purchase verb to 4.10× on a synthetic
arithmetic kernel), a cold compile cost (281 µs for the eShop purchase
verb) that recovers itself after roughly 156 invocations, and a journal
in which 99.8\% of parametric entries are compact action references
producing a footprint several-fold denser than the equivalent
literal-script storage. The magnitudes are floors of the mechanism on
RDBMS-anchored DDD aggregates whose authors designed against ORM
hydration; the same mechanism applied to host domains without that
representational pressure compounds further.

The principle takes its place alongside the analysis of the prior paper
of this series, \href{01-anti-porosity.md}{\emph{Anti-porous
Architecture}}, which characterized porosity as a representational
sparsity problem and density preservation as the operational consequence
of an event-sourced DSL journal. Where the prior paper named the defect
that domain representations on tabular substrates accumulate, this paper
names the condition under which that defect can be avoided.

{\def\LTcaptype{none} % do not increment counter
\begin{longtable}[]{@{}ll@{}}
\toprule\noalign{}
Prior paper & This paper \\
\midrule\noalign{}
\endhead
\bottomrule\noalign{}
\endlastfoot
Problem: porosity & Condition for avoiding it \\
Density preserved & Program--value separability \\
Homoiconic journal & Stable program identifier \\
Operations vs state & Functions vs instances \\
\end{longtable}
}

The two papers describe the same phenomenon from opposite sides: density
preservation is what is achieved; separability is what makes it
achievable. Together, they argue that density in domain representation
is not an optimization but a structural property that follows from how
programs relate to their values.

The principle is not bound to the runtime characterized in §3. In an
Orleans- or Akka-style actor framework with a parameterized command DSL
above a host-language domain library, the same precondition would
predict the same four-fold coherence; in a service that uses prepared
statements over an event-sourced log of named operations, three of the
four faces are already available, and the fourth --- replication-bounded
entropy --- follows once the persisted artifact is the operation rather
than its row-level effects. The realization in §3 is one example of a
construction the principle admits, not the only path to it.

Two extensions of the principle remain to be developed in subsequent
work. The fourfold coherence described here does not exhaust the
operational consequences of an externalized-parameter, locality-bound
substrate: a persistent local write buffer with asynchronous remote
replication, for instance, is structurally enabled by the same
commitments --- the actor's isolation guarantees that locally-buffered,
not-yet-replicated entries are invisible to other contexts by
construction, not by convention. The persistent buffer and its
zero-downtime implications are treated in a companion paper. A second
companion paper takes up the reaction mechanism whose
semantic-pattern-matching has been mentioned in passing throughout this
argument, and the consistency contract under which observable state
remains coherent across actors. In each case the same structural
observation continues to apply --- what makes the journal compact, what
makes deep domains compose without friction, is the separation of the
program from its values. Porosity makes deep domains expensive to
represent; anti-porosity makes them economical to compose.

\begin{center}\rule{0.5\linewidth}{0.5pt}\end{center}

\subsection{Code provenance}\label{code-provenance}

Source-code references in this paper resolve against the public
Puppeteer repository at commit
\href{https://github.com/alvaroNCubo/puppeteer/tree/2f31f9674a5de816bdf1bf9d8360ff218a02e4da}{\texttt{2f31f96}}
(2026-05-18). The snapshot is archived in Software Heritage under the
following persistent identifier:

\begin{verbatim}
swh:1:dir:10e7e6bad7eb77b6c2e406762026177f95c3ae92;
  origin=https://github.com/alvaroNCubo/puppeteer;
  anchor=swh:1:rev:2f31f9674a5de816bdf1bf9d8360ff218a02e4da
\end{verbatim}

Inline references of the form \texttt{file.cs:NN} (e.g.,
\texttt{ActorHandler.cs:38}) resolve against this snapshot. A reader can
construct a per-file SWHID by adding the qualifiers
\texttt{;path=\textless{}path\textgreater{};lines=\textless{}NN\textgreater{}}
to the directory SWHID above. Future commits to the repository may
renumber lines; the SWHID preserves the cited state independently of any
future change to the repository or its hosting.

\subsection{Acknowledgments}\label{acknowledgments}

The author used large language models (including Claude and ChatGPT) as
editorial assistants for language refinement, structural feedback, and
literature navigation. All original ideas, terminology, theoretical
constructs, and technical content presented in this work are solely the
author's.

\begin{center}\rule{0.5\linewidth}{0.5pt}\end{center}

\subsection{References}\label{references}

Appel, A. W. (1992). \emph{Compiling with continuations}. Cambridge
University Press.

Futamura, Y. (1999). Partial evaluation of computation process --- an
approach to a compiler-compiler. \emph{Higher-Order and Symbolic
Computation}, \emph{12}(4), 381--391. (Original work published 1971 in
\emph{Systems, Computers, Controls}, \emph{2}(5), 45--50.)

Hellerstein, J. M., Stonebraker, M., \& Hamilton, J. (2007).
Architecture of a database system. \emph{Foundations and Trends in
Databases}, \emph{1}(2), 141--259.

Hevner, A. R., March, S. T., Park, J., \& Ram, S. (2004). Design science
in information systems research. \emph{MIS Quarterly}, \emph{28}(1),
75--105.

International Organization for Standardization. (2023). \emph{ISO/IEC
9075:2023 Information technology --- Database languages --- SQL}.

Johnsson, T. (1985). Lambda lifting: Transforming programs to recursive
equations. In J.-P. Jouannaud (Ed.), \emph{Functional programming
languages and computer architecture} (pp.~190--203). Springer-Verlag.
(Lecture Notes in Computer Science, Vol. 201)

Jones, N. D., Gomard, C. K., \& Sestoft, P. (1993). \emph{Partial
evaluation and automatic program generation}. Prentice Hall.

Rivera, A. (2026). \emph{Anti-porous architecture: A unified design
principle for CQRS + Actor + Event-Sourcing systems} {[}Paper 1 of this
series{]}.
https://github.com/alvaroNCubo/puppeteer-papers/blob/main/01-anti-porosity.md

Stroustrup, B. (1994). \emph{The design and evolution of C++}.
Addison-Wesley.

Vandevoorde, D., Josuttis, N. M., \& Gregor, D. (2017). \emph{C++
templates: The complete guide} (2nd ed.). Addison-Wesley.

\begin{center}\rule{0.5\linewidth}{0.5pt}\end{center}

\subsection{Appendix A --- Code
references}\label{appendix-a-code-references}

The references below cite source locations in the Puppeteer codebase.
Citations are given as \texttt{file:line} pairs against the runtime
version current at the time of measurement (see §5.1 for measurement
methodology). The codebase currently lives in a private repository; a
public-mirror URL will be added to this appendix when the codebase is
sanitized for public release. Datasets cited in §5 are stored in the
companion \texttt{data/} directory of this paper repository, with the
same publication caveat.

\subsubsection{§1.2 --- Formal
characterization}\label{formal-characterization}

{\def\LTcaptype{none} % do not increment counter
\begin{longtable}[]{@{}
  >{\raggedright\arraybackslash}p{(\linewidth - 4\tabcolsep) * \real{0.3333}}
  >{\raggedright\arraybackslash}p{(\linewidth - 4\tabcolsep) * \real{0.3333}}
  >{\raggedright\arraybackslash}p{(\linewidth - 4\tabcolsep) * \real{0.3333}}@{}}
\toprule\noalign{}
\begin{minipage}[b]{\linewidth}\raggedright
Reference
\end{minipage} & \begin{minipage}[b]{\linewidth}\raggedright
Location
\end{minipage} & \begin{minipage}[b]{\linewidth}\raggedright
What it shows
\end{minipage} \\
\midrule\noalign{}
\endhead
\bottomrule\noalign{}
\endlastfoot
Runtime self-documentation as \texttt{F(x1,...,xn)} &
\texttt{ActorHandler.cs:1085} & Comment articulating the parametric
program model \\
Parametric/literal regime documented in code &
\texttt{ActorHandler.cs:1091-1093} & Comment encoding the binary
taxonomy: scripts without user parameters are interpreted, uncached,
persisted as Script entries; scripts with user parameters are compiled,
cached with an ActionId, persisted as Action entries \\
\end{longtable}
}

\subsubsection{§3.1 --- Compilation
pipeline}\label{compilation-pipeline-1}

{\def\LTcaptype{none} % do not increment counter
\begin{longtable}[]{@{}
  >{\raggedright\arraybackslash}p{(\linewidth - 4\tabcolsep) * \real{0.3333}}
  >{\raggedright\arraybackslash}p{(\linewidth - 4\tabcolsep) * \real{0.3333}}
  >{\raggedright\arraybackslash}p{(\linewidth - 4\tabcolsep) * \real{0.3333}}@{}}
\toprule\noalign{}
\begin{minipage}[b]{\linewidth}\raggedright
Reference
\end{minipage} & \begin{minipage}[b]{\linewidth}\raggedright
Location
\end{minipage} & \begin{minipage}[b]{\linewidth}\raggedright
What it shows
\end{minipage} \\
\midrule\noalign{}
\endhead
\bottomrule\noalign{}
\endlastfoot
Orchestration on first encounter & \texttt{ActorHandler.cs:1097-1148} &
\texttt{PrepareCommandProgram} rents a parser, parses the script,
decides parametric vs literal path \\
AST traversal & \texttt{Program.cs:182-244} & \texttt{ProgramExpression}
walks \texttt{Statement}s, accumulates \texttt{Expression} nodes into a
\texttt{BlockExpression} \\
Lambda composition & \texttt{Program.cs:244} &
\texttt{Expression.Lambda\textless{}Func\textless{}Parameters,\ Output,\ string\textgreater{}\textgreater{}}
wraps the \texttt{BlockExpression} \\
Compilation step & \texttt{Program.cs:163} &
\texttt{\_executable\ =\ programExpression.Compile()} \\
Eval parameter cache structure & \texttt{Program.cs:305} & Dictionary
keyed by parameter name to \texttt{(EvalScript,\ Compiled)} tuple \\
Eval recompile on script change & \texttt{Program.cs:323-331} & Cache
lookup, invalidation, re-parse, re-compile \\
\end{longtable}
}

\subsubsection{§3.2 --- Interpretation
retained}\label{interpretation-retained-1}

{\def\LTcaptype{none} % do not increment counter
\begin{longtable}[]{@{}
  >{\raggedright\arraybackslash}p{(\linewidth - 4\tabcolsep) * \real{0.3333}}
  >{\raggedright\arraybackslash}p{(\linewidth - 4\tabcolsep) * \real{0.3333}}
  >{\raggedright\arraybackslash}p{(\linewidth - 4\tabcolsep) * \real{0.3333}}@{}}
\toprule\noalign{}
\begin{minipage}[b]{\linewidth}\raggedright
Reference
\end{minipage} & \begin{minipage}[b]{\linewidth}\raggedright
Location
\end{minipage} & \begin{minipage}[b]{\linewidth}\raggedright
What it shows
\end{minipage} \\
\midrule\noalign{}
\endhead
\bottomrule\noalign{}
\endlastfoot
\texttt{CompilationModePolicy} enum & \texttt{Actor.cs:8-13} & Three
policy values (\texttt{Automatic}, \texttt{AlwaysCompiled},
\texttt{AlwaysInterpreted}); default is \texttt{Automatic} \\
Policy hint passed at call site & \texttt{ActorHandler.cs:1117-1131} &
\texttt{AdjustCompilationMode(useInterpretedMode,\ policy)} ---
\texttt{useInterpretedMode:\ true} for literal path, \texttt{false} for
parametric \\
Policy resolution switch & \texttt{Program.cs:134-150} & Sets
\texttt{IsCompiledMode} from policy + hint \\
Execution dispatch & \texttt{ActorHandler.cs:1955-1968} &
\texttt{Perform} calls \texttt{ExecuteExpression} if compiled,
\texttt{Execute} otherwise \\
Commands caching rule & \texttt{ActorHandler.cs:1117} & Caches only when
\texttt{parameters.HasUserParameter()} is true \\
Queries caching rule & \texttt{ActorHandler.cs:1405} & Caches when
\texttt{EMPTY\_PARAMETERS} \emph{or} \texttt{HasUserParameter()} is true
(asymmetry vs commands) \\
\end{longtable}
}

\subsubsection{§3.3 --- Hot-loaded DSL
programs}\label{hot-loaded-dsl-programs}

{\def\LTcaptype{none} % do not increment counter
\begin{longtable}[]{@{}
  >{\raggedright\arraybackslash}p{(\linewidth - 4\tabcolsep) * \real{0.3333}}
  >{\raggedright\arraybackslash}p{(\linewidth - 4\tabcolsep) * \real{0.3333}}
  >{\raggedright\arraybackslash}p{(\linewidth - 4\tabcolsep) * \real{0.3333}}@{}}
\toprule\noalign{}
\begin{minipage}[b]{\linewidth}\raggedright
Reference
\end{minipage} & \begin{minipage}[b]{\linewidth}\raggedright
Location
\end{minipage} & \begin{minipage}[b]{\linewidth}\raggedright
What it shows
\end{minipage} \\
\midrule\noalign{}
\endhead
\bottomrule\noalign{}
\endlastfoot
Domain library loading & \texttt{DomainLibraries.cs:77-115} &
\texttt{GetOrLoad(params\ Assembly{[}{]})} reflects public types and
caches them once per deduplicated assembly set; a single-assembly
overload remains for the back-compat path \\
Library binding to actor & \texttt{ActorHandler.cs:61} &
\texttt{libraries\ =\ DomainLibraries.GetOrLoad(LibraryAssemblies)} \\
Fluent script invocation & \texttt{ActorV2.cs:32} &
\texttt{Using(scriptForChk,\ scriptForCmd)} introduces new scripts at
runtime against the cached domain \\
\end{longtable}
}

\subsubsection{§4 --- Verb richness}\label{verb-richness-1}

{\def\LTcaptype{none} % do not increment counter
\begin{longtable}[]{@{}
  >{\raggedright\arraybackslash}p{(\linewidth - 4\tabcolsep) * \real{0.3333}}
  >{\raggedright\arraybackslash}p{(\linewidth - 4\tabcolsep) * \real{0.3333}}
  >{\raggedright\arraybackslash}p{(\linewidth - 4\tabcolsep) * \real{0.3333}}@{}}
\toprule\noalign{}
\begin{minipage}[b]{\linewidth}\raggedright
Reference
\end{minipage} & \begin{minipage}[b]{\linewidth}\raggedright
Location
\end{minipage} & \begin{minipage}[b]{\linewidth}\raggedright
What it shows
\end{minipage} \\
\midrule\noalign{}
\endhead
\bottomrule\noalign{}
\endlastfoot
Check-then-command fluent dispatch & \texttt{ActorV2Invocation.cs:69} &
\texttt{PerformCheckThenCommand()} --- read-lock check + write-lock
command + journal write, with rollback on check failure \\
\end{longtable}
}

\subsubsection{§5 --- Empirical datasets}\label{empirical-datasets}

The datasets that produced the magnitudes reported in §5 are stored in
the companion \texttt{data/} directory of the paper repository. The
production-verb measurements are produced against two independent
open-source aggregates:

\begin{itemize}
\tightlist
\item
  \textbf{eShop} ---
  \href{https://github.com/dotnet/eShop}{\texttt{dotnet/eShop}} (MIT),
  \texttt{src/Ordering.Domain/AggregatesModel/OrderAggregate/Order.cs}.
  Multi-item purchase via 10-arg constructor plus four
  \texttt{AddOrderItem} calls plus a four-step state-machine walk to
  Shipped.
\item
  \textbf{Grzybek} ---
  \href{https://github.com/kgrzybek/modular-monolith-with-ddd}{\texttt{kgrzybek/modular-monolith-with-ddd}}
  (MIT),
  \texttt{src/Modules/Payments/Domain/SubscriptionPayments/SubscriptionPayment.cs}.
  Subscription payment lifecycle via \texttt{Buy} factory plus
  \texttt{MarkAsPaid}.
\end{itemize}

The synthetic kernels in §5.2 (arithmetic, DSL-rich) carry no
external-codebase dependency.

{\def\LTcaptype{none} % do not increment counter
\begin{longtable}[]{@{}
  >{\raggedright\arraybackslash}p{(\linewidth - 8\tabcolsep) * \real{0.2000}}
  >{\raggedright\arraybackslash}p{(\linewidth - 8\tabcolsep) * \real{0.2000}}
  >{\raggedright\arraybackslash}p{(\linewidth - 8\tabcolsep) * \real{0.2000}}
  >{\raggedright\arraybackslash}p{(\linewidth - 8\tabcolsep) * \real{0.2000}}
  >{\raggedright\arraybackslash}p{(\linewidth - 8\tabcolsep) * \real{0.2000}}@{}}
\toprule\noalign{}
\begin{minipage}[b]{\linewidth}\raggedright
Section
\end{minipage} & \begin{minipage}[b]{\linewidth}\raggedright
Dataset directory
\end{minipage} & \begin{minipage}[b]{\linewidth}\raggedright
Lab
\end{minipage} & \begin{minipage}[b]{\linewidth}\raggedright
Host
\end{minipage} & \begin{minipage}[b]{\linewidth}\raggedright
What it contains
\end{minipage} \\
\midrule\noalign{}
\endhead
\bottomrule\noalign{}
\endlastfoot
§5.2 (compilation speedup, production verb) & \texttt{data/lab01-eshop/}
& Lab 1 Run 3 & eShop & Compiled vs interpreted bench, N=1000,
bootstrap-then-measurement \\
§5.2 (compile cold cost, production verb) & \texttt{data/lab02-eshop/} &
Lab 2 Tier 3 & eShop & Cold compile cost per distinct script variant,
N=100 \\
§5.3 (eval cache hit vs miss) & \texttt{data/lab03-eshop/} & Lab 3 Tier
C & eShop & Stable vs mutating eval text; per-iteration end-to-end and
isolated eval-compile ticks \\
§5.4 (journal density) & \texttt{data/lab04-eshop/} & Lab 4 & eShop &
Action / Script / Define entry counts and payload bytes parsed directly
from BinaryEventCodec journal\_*.bin \\
§5.4 (journal density, replication) & \texttt{data/lab04-grzybek/} & Lab
4 & Grzybek & Same entry-type analysis applied to the
SubscriptionPayment lifecycle verb \\
§5.5 (verb richness) & \texttt{data/lab05-eshop/} & Lab 5 & eShop & α
DSL dispatch counter (runtime) plus β Roslyn forward closure over
\texttt{Ordering.Domain} \\
§5.5 (verb richness, replication) & \texttt{data/lab05-grzybek/} & Lab 5
& Grzybek & α + β over \texttt{Payments.Domain} plus shared
\texttt{BuildingBlocks/Domain} \\
\end{longtable}
}

Each dataset directory contains a \texttt{headline.md} summary, raw
CSVs, and a Git SHA stamping the runtime version against which the lab
was run. The Roslyn walker source for the β half of §5.5 lives at
\texttt{labs/lab05-eshop-roslyn/} and
\texttt{labs/lab05-grzybek-roslyn/} --- both are self-contained .NET 9
console projects that can be re-executed against the public source trees
of the respective aggregates.

\end{document}
